\chapter{Rapid PVST+ and the Protection Features}\label{ch:stp}
\bp{2.5}

\begin{outcomes}
By the end of this chapter you will be able to:
\begin{tight}
  \item \textbf{Explain} what a layer 2 loop does to a network, and why it is
        worse than any layer 3 failure.
  \item \textbf{Predict} the root bridge, root ports and designated ports for a
        given topology, and \textbf{influence} the outcome deliberately.
  \item \textbf{Configure} Rapid PVST+, root primary and secondary, PortFast,
        BPDU guard, root guard and loop guard.
  \item \textbf{Diagnose} an unexpected root bridge, an errdisabled port and a
        port that will not leave blocking.
\end{tight}
\end{outcomes}

\begin{prereq}
\chref{ethernet} for flooding --- which is what makes a loop catastrophic ---
and \chref{trunking} for trunks between switches.
\end{prereq}

\begin{keyterms}
BPDU $\bullet$ root bridge $\bullet$ bridge ID $\bullet$ bridge priority
$\bullet$ root port $\bullet$ designated port $\bullet$ alternate port $\bullet$
path cost $\bullet$ Rapid PVST+ $\bullet$ PortFast $\bullet$ BPDU guard
$\bullet$ root guard $\bullet$ loop guard $\bullet$ errdisable
\end{keyterms}

\section{Why a layer 2 loop is fatal}

\begin{conceptbox}[Three things happen at once, and none of them stops]
An IP packet carries a TTL that decrements at each router, so a routing loop
eventually discards the packet. \textbf{An Ethernet frame has no TTL.} A frame
caught in a switching loop circulates until somebody unplugs something.

\begin{tight}
  \item \textbf{Broadcast storm.} A single broadcast is flooded, arrives back,
        and is flooded again --- multiplying without limit until the links
        saturate.
  \item \textbf{MAC table instability.} The same source address arrives on two
        ports repeatedly, so the table is rewritten constantly and unicast
        forwarding becomes unreliable (see the flap message in \chref{ethernet}).
  \item \textbf{Duplicate frame delivery.} Hosts receive multiple copies, which
        breaks protocols that assume they will not.
\end{tight}
The network does not degrade gracefully; it stops, usually within seconds. This
is why spanning tree exists and why it is on by default.
\end{conceptbox}

\ccnafig{%
\begin{tikzpicture}[font=\scriptsize]
  \tikzset{
    dv/.style={rounded corners=3pt, draw=#1, fill=#1!8, text=#1!45!black,
               line width=0.9pt, minimum width=1.4cm, minimum height=0.95cm,
               align=center, font=\tiny\bfseries},
    w/.style={line width=1pt, neutral!70},
    hdr/.style={font=\scriptsize\bfseries, text=neutral!45!black}}

  %======== left: the loop
  \node[hdr] at (2.5,2.95) {Without spanning tree};
  \node[dv=dom2] (a1) at (2.5,2.05) {\faIcon{sitemap}\\SW1};
  \node[dv=dom2] (b1) at (1.1,0.25) {\faIcon{sitemap}\\SW2};
  \node[dv=dom2] (c1) at (3.9,0.25) {\faIcon{sitemap}\\SW3};
  \draw[w] (a1) -- (b1);  \draw[w] (a1) -- (c1);  \draw[w] (b1) -- (c1);

  % Sits inside the triangle, clear of all three switch boxes. An arc rather
  % than a circle: TikZ puts no arrowhead on a closed path, so a circle drawn
  % with an arrow tip renders as a plain ring with no sense of direction.
  \draw[-{Stealth[length=2.6mm]}, accent!85, line width=1.6pt]
      (2.96,0.95) arc (0:308:0.46);
  \node[font=\tiny\bfseries, text=accent!75!black, align=center] at (2.5,-0.95)
      {one broadcast circulates\\for ever};

  %======== right: the tree
  \node[hdr] at (10.5,2.95) {With spanning tree};
  \node[dv=dom2] (a2) at (10.5,2.05) {\faIcon{sitemap}\\SW1};
  \node[font=\tiny\bfseries, text=dom2!50!black] at (10.5,3.0) {};
  \node[dv=dom2] (b2) at (9.1,0.25) {\faIcon{sitemap}\\SW2};
  \node[dv=dom2] (c2) at (11.9,0.25) {\faIcon{sitemap}\\SW3};
  \draw[w] (a2) -- (b2);  \draw[w] (a2) -- (c2);
  \draw[w, accent!70, dash pattern=on 3pt off 2.5pt] (b2) -- (c2);

  \node[circle, draw=accent!85, fill=white, line width=1.1pt,
        inner sep=1.2pt, font=\tiny\bfseries, text=accent!75!black]
      at (10.5,0.25) {BLK};
  \node[font=\tiny\bfseries, text=dom2!45!black, anchor=west] at (11.3,2.05)
      {ROOT};
  \node[font=\tiny, text=neutral!85!black, align=center] at (10.5,-0.95)
      {one port blocks, so there is\\exactly one path between any two};

  \node[font=\scriptsize\bfseries, text=neutral!45!black] at (6.5,1.2)
      {$\Longrightarrow$};
\end{tikzpicture}}%
{Spanning tree does not remove the cable. It blocks one port, so the physical
redundancy is still there and can be brought back automatically when a link
fails.}{stploop}

\section{How the tree is chosen}

\begin{definitionbox}[Bridge ID and root election]
Every switch has a \term{bridge ID}: a 4-bit priority, plus a 12-bit extended
system ID carrying the VLAN number, plus its MAC address. The switch with the
\textbf{lowest bridge ID wins} and becomes the \term{root bridge}.

Default priority is 32768. Because the extended system ID adds the VLAN number,
the effective priority in VLAN 10 is $32768+10=32778$. Priorities are configured
in steps of 4096.
\end{definitionbox}

\begin{worked}[Which of these three is the root, in VLAN 10?]
Three switches, all left at defaults:

\begin{center}
\begin{tabular}{@{}lll@{}}
\cmd{SW1} & priority 32768 & MAC \cmd{0011.2233.0100} \\
\cmd{SW2} & priority 32768 & MAC \cmd{0011.2233.0050} \\
\cmd{SW3} & priority 32768 & MAC \cmd{0011.2233.0200} \\
\end{tabular}
\end{center}

\textbf{Step 1 --- compare priority.} All three are 32768, and all three are in
VLAN 10, so all three have the same effective priority of $32768+10 = 32778$.
No winner yet.

\textbf{Step 2 --- the tie-break is the MAC address, and lowest wins.} Compare
them left to right, as you would compare words alphabetically. The first eight
digits are identical, so the decision comes down to the last four:
\cmd{0100}, \cmd{0050}, \cmd{0200}.

In hex, \cmd{0050} is the smallest of the three (\chref{ethernet} if the hex is
not yet automatic --- \cmd{0050} is 80 in decimal, \cmd{0100} is 256).

\textbf{\cmd{SW2} is the root bridge.}

\medskip
Now notice what just happened: \textbf{nobody chose this}. The winner is
whichever switch happens to have the lowest MAC address, which correlates with
nothing except manufacturing date. That is the point of the alert below.
\end{worked}

\begin{alertbox}[Left alone, the oldest switch wins]
If every switch keeps the default priority, the election is decided by the lowest
MAC address --- which usually means the oldest device in the building becomes the
root of your network. That is almost never where you want traffic to converge.
\textbf{Always set the root deliberately.}
\end{alertbox}

\begin{conceptbox}[Then every port takes a role]
Once the root is known:
\begin{tight}
  \item Each non-root switch picks its \term{root port} --- the one port with the
        lowest cumulative path cost to the root. It forwards.
  \item Each segment picks a \term{designated port} --- the port on that segment
        closest to the root. It forwards.
  \item Everything else becomes an \term{alternate} (or backup) port and
        \textbf{blocks}. That is the loop being broken.
\end{tight}
Ties are broken by lowest sender bridge ID, then lowest sender port ID.
\end{conceptbox}

\begin{longtable}{@{}L{4.4cm}C{3.0cm}L{7.6cm}@{}}
\toprule
\rowcolor{primary!15}
\textbf{Link speed} & \textbf{Cost} & \textbf{Note} \\
\midrule
\endhead
10 Mbps   & 100 & Legacy \\
\rowcolor{lightbg}
100 Mbps  & 19  & \\
1 Gbps    & 4   & \\
\rowcolor{lightbg}
10 Gbps   & 2   & Costs are cumulative along the path to the root \\
\bottomrule
\end{longtable}

\section{Configuring Rapid PVST+}

\term{Rapid PVST+} runs one rapid spanning tree instance per VLAN. It converges
in seconds rather than the 30--50 of the original standard, and it is what the
blueprint names.

\begin{config}{Rapid PVST+ with a deliberate root}
! Run rapid per-VLAN spanning tree.
spanning-tree mode rapid-pvst
!
! On the switch that SHOULD be root for these VLANs:
spanning-tree vlan 10,20,30 root primary
!
! On the switch that should take over if it fails:
!   spanning-tree vlan 10,20,30 root secondary
!
! The macro above simply sets a priority. You can be explicit:
!   spanning-tree vlan 10 priority 4096
\end{config}

\begin{verify}{SW1}{show spanning-tree vlan 10}
VLAN0010
  Spanning tree enabled protocol rstp
  Root ID    Priority    4106
             Address     00d0.58a1.0200
             This bridge is the root
             Hello Time  2 sec  Max Age 20 sec  Forward Delay 15 sec

  Bridge ID  Priority    4106   (priority 4096 sys-id-ext 10)
             Address     00d0.58a1.0200

Interface           Role Sts Cost      Prio.Nbr Type
------------------- ---- --- --------- -------- --------------------------------
Gi0/1               Desg FWD 4         128.1    P2p
Gi0/2               Desg FWD 4         128.2    P2p
\end{verify}

Two things to read first: \cmd{This bridge is the root} --- is it the switch you
intended? --- and the priority, \cmd{4106}, which is $4096+10$ and confirms the
extended system ID is doing what it should.

\begin{verify}{SW2}{show spanning-tree vlan 10}
VLAN0010
  Spanning tree enabled protocol rstp
  Root ID    Priority    4106
             Address     00d0.58a1.0200
             Cost        4
             Port        1 (GigabitEthernet0/1)

Interface           Role Sts Cost      Prio.Nbr Type
------------------- ---- --- --------- -------- --------------------------------
Gi0/1               Root FWD 4         128.1    P2p
Gi0/2               Altn BLK 4         128.2    P2p
\end{verify}

\cmd{Gi0/2} is \cmd{Altn BLK} --- alternate, blocking. That is the loop being
broken, and it is a healthy state, not a fault.

\section{The protection features}

Topic \tp{2.5} names PortFast, root guard, loop guard and BPDU guard explicitly.

\begin{longtable}{@{}L{2.8cm}L{5.6cm}L{6.6cm}@{}}
\toprule
\rowcolor{primary!15}
\textbf{Feature} & \textbf{What it does} & \textbf{Put it where} \\
\midrule
\endhead
\textbf{PortFast} & Skips listening and learning; the port forwards immediately &
Access ports with an end device. Saves ~30 s on boot and stops DHCP timing out \\
\rowcolor{lightbg}
\textbf{BPDU guard} & Errdisables the port if a BPDU arrives &
The same access ports. A BPDU there means somebody plugged in a switch \\
\textbf{Root guard} & Ignores superior BPDUs; puts the port in root-inconsistent
state & Ports facing downstream switches you do not want becoming root \\
\rowcolor{lightbg}
\textbf{Loop guard} & Blocks a port that stops receiving BPDUs rather than
letting it forward & Root and alternate ports on trunks, where a unidirectional
link would otherwise create a loop \\
\bottomrule
\end{longtable}

\begin{alertbox}[PortFast and BPDU guard belong together]
PortFast on its own is a loop waiting to happen: it tells a port to forward
immediately, so if somebody plugs a switch into it the loop is live before
spanning tree can react. BPDU guard is the safety catch --- the instant a BPDU
arrives, the port is shut. Configure them as a pair, always.
\end{alertbox}

\begin{config}{protection, per port and globally}
! Per access port.
interface range GigabitEthernet0/5 - 20
 switchport mode access
 switchport access vlan 10
 spanning-tree portfast
 spanning-tree bpduguard enable
!
! Or globally, which applies to every access port automatically --
! far harder to forget on the next port somebody configures.
spanning-tree portfast default
spanning-tree portfast bpduguard default
!
! On a trunk facing a downstream switch that must never be root:
interface GigabitEthernet0/23
 spanning-tree guard root
!
! Loop guard, globally on all point-to-point links.
spanning-tree loopguard default
\end{config}

\begin{verify}{SW1}{show interfaces status err-disabled}
Port      Name               Status       Reason               Err-disabled Vlans
Gi0/7                        err-disabled bpduguard
\end{verify}

\begin{conceptbox}[Recovering an errdisabled port]
Manually: \cmd{shutdown} then \cmd{no shutdown} on the interface --- after
removing whatever caused it. Automatically:

\begin{config}{automatic recovery from BPDU guard}
errdisable recovery cause bpduguard
errdisable recovery interval 300
\end{config}

Automatic recovery is a convenience, not a fix. If a rogue switch is still
plugged in the port will simply errdisable again every five minutes.
\end{conceptbox}

\begin{pitfall}
\begin{tight}
  \item \textbf{Leaving the root election to chance.} The oldest switch wins, and
        your traffic detours through the wiring closet nobody has visited in ten
        years.
  \item \textbf{PortFast without BPDU guard.} See above.
  \item \textbf{PortFast on a trunk to another switch.} It is for edge ports.
        There is a \cmd{portfast trunk} variant for specific cases, but the
        default answer is no.
  \item \textbf{Treating a blocking port as broken.} \cmd{Altn BLK} is spanning
        tree working.
  \item \textbf{Forgetting spanning tree is per VLAN.} A change to VLAN 10 does
        nothing for VLAN 20. Always name the VLANs.
\end{tight}
\end{pitfall}

\begin{breakfix}[After a new switch was added, traffic takes a strange path.]
An access switch was installed in a spare closet. Since then, throughput between
two core switches has halved and \cmd{show spanning-tree vlan 10} on the core
reports that it is no longer the root.

\begin{verify}{CORE1}{show spanning-tree vlan 10 | include Root ID|Priority|Address}
  Root ID    Priority    32778
             Address     0011.2233.4455
             Cost        19
\end{verify}

\textbf{Diagnosis.} The root's priority is 32778 --- the default of $32768+10$ ---
and the MAC address is not the core's. The new access switch shipped with default
priority and a lower MAC address than the core, so it won the election. Every
VLAN's traffic now converges on a small access switch connected by a 100~Mbps
link, which is where the cost of 19 comes from and why throughput halved.

\textbf{Fix.} Two steps, and do both. Set the intended root explicitly on the
core: \cmd{spanning-tree vlan 10,20,30 root primary}, and the secondary on the
other core switch. Then stop it happening again: apply
\cmd{spanning-tree guard root} on the core ports facing access switches, so a
switch advertising a superior BPDU is ignored rather than obeyed.
\end{breakfix}

\begin{lab}{Steer the tree, then protect it}{Packet Tracer}
\textbf{Topology.} Three switches cabled in a triangle --- deliberately looped ---
with PCs on two of them.

\begin{tasks}
  \item Cable the triangle with no spanning tree changes. Identify the root from
        \cmd{show spanning-tree} and explain \emph{why} that switch won.
  \item Record which port on which switch is blocking.
  \item Set \cmd{SW1} as \cmd{root primary} and \cmd{SW2} as \cmd{root secondary}
        for VLAN 1. Re-check and record what moved.
  \item Disconnect the link between \cmd{SW1} and \cmd{SW3}. Time how long
        traffic takes to recover and note which port changed role.
  \item Enable PortFast and BPDU guard on an access port, then plug a switch into
        it. Record the log message and the port status, then recover the port.
  \item Enable \cmd{spanning-tree guard root} on \cmd{SW1}'s port to \cmd{SW3},
        then try to make \cmd{SW3} root by lowering its priority. Record what
        happens instead.
  \item \textbf{Extension.} Time how long the network takes to recover from a
        link failure under \cmd{spanning-tree mode pvst}, then repeat under
        \cmd{rapid-pvst} and compare the two figures.
\end{tasks}
\end{lab}

\begin{checkpoint}
\begin{tightnum}
  \item Three switches have default priorities and MAC addresses ending
        \cmd{...0100}, \cmd{...0050} and \cmd{...0200}. Which is root, and why?
  \item A port shows \cmd{Altn BLK}. Is this a fault? Justify your answer.
  \item Why must PortFast and BPDU guard be configured together?
\end{tightnum}
\end{checkpoint}

\begin{chaptersummary}
\begin{tight}
  \item An Ethernet frame has no TTL, so a layer 2 loop never self-clears. A
        broadcast storm takes the network down in seconds.
  \item Lowest bridge ID becomes root; bridge ID is priority plus VLAN plus MAC.
        Set it deliberately or the oldest switch wins.
  \item Root port per switch, designated port per segment, everything else
        blocks.
  \item Rapid PVST+ runs one rapid instance per VLAN. Configure it, and remember
        every command is per VLAN.
  \item PortFast with BPDU guard on access ports; root guard facing downstream
        switches; loop guard on trunks. Errdisable is recovered by
        \cmd{shutdown}/\cmd{no shutdown} after removing the cause.
\end{tight}
\end{chaptersummary}

\begin{reviewq}
  \item \textbf{(MCQ)} What is the effective priority of a switch in VLAN 20 with
        a configured priority of 8192? (a)~8192 (b)~8212 (c)~28692 (d)~32788
  \item \textbf{(MCQ)} Which feature errdisables a port when a BPDU is received?
        (a)~PortFast (b)~BPDU guard (c)~root guard (d)~loop guard
  \item \textbf{(Short)} Explain why a layer 2 loop is more damaging than a
        routing loop.
  \item \textbf{(Short)} A port is in \cmd{root-inconsistent} state. Which feature
        put it there, and what was it protecting against?
  \item \textbf{(Applied)} You inherit a three-switch network where the root is a
        random access switch and no protection features are configured. Write the
        commands to make the core the root, provide a backup root, and prevent
        recurrence --- and say in what order you would apply them on a live
        network, with a reason.
\end{reviewq}
